\chapter{Discussion and Conclusions}\label{ch:discussion}

\section{What the results say}\label{sec:disc-findings}

Most of the variation in K16 is extrinsic, in the same sense used by Elowitz et al.\ \cite{elowitz2002}. Cell-to-cell differences in brightness throughout the movie are much larger than randomness specific to the promoter regions. Hence, the switching rates are poorly constrained. It, again, as seen in Elowitz, would also suggest that a two-color reporter, which further separates intrinsic from extrinsic noise directly, would help more than simply adding cells. That is to say, the model describes the population but not the individual cells it tries to model in the first place. Because it reproduces every summary statistic, the model expects them to revert to some mean within a lifetime, while the actual cells would steadily drift and fluctuate over time. Moreover, when considering that I am explicitly looking at Tat, which is known to have a positive feedback loop, the model ends up ignoring known cell-cycle effects. The missing mRNA step explains why the ON fraction is only an effective quantity, while burst frequency is still comparable with the lab's estimate.

\section{Limitations and deviations from the plan}\label{sec:disc-deviations}

Only one well was analyzed, and it is not untreated. K16 received compound NSC 157725 (lab ID D62), which Lu et al.\ \cite{lu2021} list as an HF-$\tau_{1/2}$ ``Decrease'' hit (2.35\,h against a median of 3.16\,h over 34 plate-2 control wells), so the circuit here is perturbed. The background level is uncertain, drift removal shifts the ON fraction, and burst size is only in arbitrary units. Cell division is not modeled either: one daughter keeps the track and the other is dropped.

There were three deviations from the plan. One well was analyzed instead of $\sim$1,800 conditions, because of the data size and segmentation time. ABC replaced particle MCMC because one reusable reference table, built in parallel, made the calibration tests affordable. In hindsight the forecasting failure backs this up, since a full trajectory likelihood would have been dominated by the model's misfit to individual cells. The particle filter was still built, and used for forecasting. Finally, the feedback model and the DNA demonstration were not done. The DNA link stays theoretical: any chemical reaction network, including the inferred telegraph model, can in principle be compiled into a DNA strand-displacement circuit \cite{soloveichik2010}. That time went into the diagnostic study, the exact-likelihood benchmark, the large-number model, the held-out tests and reproducing the paper's metrics.

\section{What I learned}\label{sec:disc-learning}

I think the most important takeaway is to treat validation as an initial part of the design rather than a single final step. While this was the initial goal, it was very useful to have a baseline at each point and know when to start asking questions or when it was safe to move on. Every stage of this project had some layer of redundancy. This was most useful with regards to the baseline ABC, where the interval happened to contain the true interval, and I otherwise might have moved on without further investigation. That takeaway is to consider calibration and width together, and check calibration with a coverage test rather than a handful of examples. On the computing side, vectorization became a practical question rather than a rule. Gillespie cannot be vectorized with respect to time, but cells and table rows are independent, and working out which axis could be vectorized decided most of the performance. I was already using EC2 for a MAST90098 project, and running the reference table there was a direct use of the parallel-computing ideas from this one. Even so, switching to the large-number model saved more time than the extra hardware did. Building the filter self-test taught me to separate ``the model is wrong'' from ``the code is wrong'' before drawing conclusions, and the forecasting failure said more about the biology than the fit did. I also found it unintuitive that burst size is a ratio, $k_m/k_{off}$, rather than a rate. That is why it can be well determined when neither rate is.

\section{Future work}\label{sec:disc-future}

The natural next step is to add an mRNA step and Tat feedback to the model, and to let the rates vary over time instead of removing the drift. Fitting the untreated and control wells from the authors' traces, which are already downloaded, would give a like-for-like comparison with the lab's model. With a validated particle filter already in place, particle MCMC is also within reach.

\section{Conclusion}\label{sec:conclusion}

The purpose of this project was to see whether I could recover hidden promoter parameters from the HIV latency circuit using single-cell fluorescence movies. Over the course of 7 weeks, it achieved each goal it set out to, excluding the additional extensions, as time was instead spent debugging and addressing issues caught by the earlier layers of redundancy. That is, I derived the traces and luminosity using the pixels from the movie, built and validated a Gillespie simulator, and finally performed and validated inference on the posterior for one well. Through this process, I found that rejection ABC lost information, and found that a regression adjustment could recover it. That well K16 is dominated by extrinsic noise, and the inferred burst parameters agree with the lab's independent estimate. Finally, the model better captures the entire population than individual cells, which are likely a result of the dynamics it leaves out.

\section*{Code availability}\label{sec:code-availability}
\addcontentsline{toc}{section}{Code availability}

The code is submitted with this report in the \texttt{FinalCode} folder. It needs Python 3.12 with the pinned versions in \texttt{requirements.txt}, and the notebooks run in order from \texttt{01\_extract\_traces.ipynb} to \texttt{05\_hiv\_inference.ipynb}. The reference tables are precomputed in \texttt{results/} and can be rebuilt with \texttt{build\_reference\_table.py}; the README gives the EC2 steps. Tests run with \texttt{python -m pytest tests}. Appendix~\ref{app:reproduce} maps every figure and table to its code, and the data are available from the Illinois Data Bank \cite{idb_dataset}.
